\section{Limitations}\label{sec:limitations}

\emph{What the guarantees assume.} Exact coverage needs exchangeable families (A1), or a family-level shift whose density ratio is a nonnegative function in the span of the features; with an estimated ratio it needs the ratio to be fitted on families independent of those it calibrates (Lemma~\ref{lem:crossfit}), and the class-conditional versions hold for families that contain the class. Every guarantee is marginal over families: ProtCOVER certifies no particular protein, and the coverage of a small family is dominated by binomial noise. Under any other shift, Corollary~\ref{cor:conditional} bounds the loss by total-variation distances that cannot be estimated from unlabeled data, and when a family's scores and labels given its descriptors differ between calibration and deployment (A2 fails, as in S5), every method under-covers and the ratio gives no warning.

\emph{What is only measured.} The comparison with the baselines, the robustness to the family definition, the base model and missing alignments, and the external results are empirical. Beyond the 0.01 margin, ProtCOVER's advantage over group-weighted conformal prediction is confined to descriptor shift (S2) on one benchmark whose ten seeds resplit one pool of families; on orphan families, held-out embedding regions and the external evaluations it matches group weighting. The advantage is in coverage, not efficiency: ProtCOVER buys its smaller shortfall with larger sets under shift.

\emph{Scope.} The ratio-based version needs a batch of unlabeled target families, or a reference panel such as the unlabeled proteome for a single query; ProtCOVER-D does not, at a somewhat larger shortfall. The base model is deliberately simple, and a stronger one may give smaller sets. ProtCOVER outputs prediction sets, not ACMG/AMP evidence strengths; mapping sets to PP3/BP4 strengths is future work.
